
%%%%%%%%%%%%%%%%%%%%%%%%%%%%%%%%%%%%%%%%%%%%%%%%%%%%%%%%%%%%%
\section{模型假设}

结合题目附录所给出的干扰源特征、测向机工作原理与模拟器行为规则，本文做出以下假设。

\begin{itemize}[itemindent=2em]
\item \textbf{假设1：}干扰源的示向度误差是随地点固定、随检测点独立的随机量。同一检测点处的电磁环境在一段时间内固定，重复检测不会改变误差，因此同一检测点的多次观测不提供额外的角度信息；不同检测点的误差相互独立，且在 $[-1^\circ,1^\circ]$ 内服从均匀分布。该假设与附录 1 的说明一致，是问题二中把示向度误差按最坏情形取为上界、在蒙特卡洛中按均匀分布抽样的依据。
\item \textbf{假设2：}干扰源在有效覆盖角度范围内信号均匀辐射，场强仅随距离衰减、与角度无关，因而"能否检测"只取决于检测点到干扰源的距离是否超过有效接收半径，以及检测点是否落在定向源的有效覆盖角度范围内。据此，全向源的可检测域为圆盘，定向源的可检测域为半圆盘。
\item \textbf{假设3：}各干扰源的信号互不影响，测向机在某个频道上检测到的示向度只由该频道对应的干扰源决定，不受其他频道上干扰源的干扰。因此不同频道的检测结果相互独立，可以逐频道独立处理。
\item \textbf{假设4：}每个干扰源的有效接收半径在 1000 米到 1500 米之间且通过接口不可获知，因此本文在"保证可测"的设计中统一取最坏情形下的下界 1000 米作为保证半径，即只要观测点与干扰源的距离不超过 1000 米，就认为该干扰源一定能够被测到。该处理是保守的，其代价是观测网略密于理论最优，收益是策略的正确性不依赖于对未知参数的猜测。
\item \textbf{假设5：}机器狗在两次检测之间的移动沿直线进行，速度恒为 $5$ 米每秒；移动过程中无法进行有效检测，必须停止移动后完成检测。该假设决定了移动时间等于前后两点间的直线距离除以速度，是虚拟时间模型中移动项的直接来源。
\item \textbf{假设6：}在目标区域的尺度（半径 1800 米）下，干扰源的位置被视为平面上的点，测向机与所有干扰源位于同一平面，忽略高度差异；光学精确定位与清除操作的成功与否仅取决于清除点与干扰源之间的距离是否不超过 20 米，与干扰源的信号覆盖角度范围无关。
\item \textbf{假设7：}干扰源个数与位置在目标区域内按均匀分布生成，即不存在事先已知的热点区域。该假设用于问题二可行域采样、问题一与问题三的蒙特卡洛统计以及理论下界的推导，当实际分布偏离均匀时本文以灵敏度分析的形式讨论其影响。
\item \textbf{假设8：}任务评价的时间口径为虚拟世界时间，即移动时间、频道切换时间、检测时间与清除时间之和；题目中 20 分钟的限制针对机器狗程序的现实运行时长，两种时间口径相互独立。该假设是问题三与问题四中"最小化时间"这一提法得以明确化的前提。
\end{itemize}
